\par\textbf{Procedimiento completo para comparar tu solución.} Con eficiencia y TAS constantes, $dW/dt=-g_0c_mW/E$. Separar $dW/W=-(g_0c_m/E)dt$ e integrar entre pesos da tiempo logarítmico. La distancia usa TAS constante y la masa consumida se obtiene dividiendo diferencia de pesos por gravedad.

\begin{enumerate}

\item Evaluar la integral temporal.

\[\fitmath{t=\frac{15}{9.80665\cdot 1.6\times10^{-5}}\cdot \ln\left(\frac{36000}{33000}\right)\approx 8318.148\,\mathrm{s}}\]

\item Multiplicar por velocidad constante.

\[\fitmath{R=100\cdot t}\]

\[\fitmath{R=100\cdot (8318.147984)\approx 831814.8\,\mathrm{m}}\]

\item Convertir peso perdido en masa consumida.

\[\fitmath{\mathrm{fuel}=\frac{36000-33000}{9.80665}\approx 305.9149\,\mathrm{kg}}\]

\end{enumerate}
